
To write $f$ in quadratic form
\[
f(x) = \frac{1}{2} x^T Q x + g^T x + c,
\]
we calculate its derivatives:
\[
\nabla f(x) = \cvect{x_1-x_2+3 \\ -x_1+x_2} \text{ and } \nabla^2 f(x)
= \begin{pmatrix*}
  1 & -1 \\ -1 & 1
\end{pmatrix*}.
\]
Therefore,
\[
Q=\begin{pmatrix*}
  1 & -1 \\ -1 & 1
\end{pmatrix*}, \; g=\cvect{3 \\ 0}, \; c= -5.
\]


The positive definiteness of $Q$ can be identified from its
eigenvalues, that are the roots of the characteristic
polynomial. Therefore, we solve the equation
\begin{align*}
\det(\lambda I - Q) = \det \begin{pmatrix}
1- \lambda & -1\\ -1 & 1-\lambda
\end{pmatrix} & = (1-\lambda)^2 - 1\\
& = 1-2\lambda+\lambda^2 -1 \\
&= \lambda(\lambda-2).
\end{align*}

Therefore, the eigenvalues of $Q$ are $\lambda_1=2$ and $\lambda_2=0$. We deduce that $Q$ is positive semidefinite. Thus, either the problem is not bounded or there is an infinite number of global minima.

In order to find in which case we are, we consider the Schur decomposition of $Q$. To do so, we first have to determine the eigenvectors associated with each eigenvalue.

\begin{itemize}
\item For $\lambda_1$ we have $$\begin{pmatrix*}[r]
  -1&-1\\-1&-1\end{pmatrix*} \begin{pmatrix*}[r] x_1
    \\ x_2\end{pmatrix*} = \begin{pmatrix*}[r] 0\\0\end{pmatrix*},$$
      which implies that the normalized eigenvector is
      \[
      x_{\lambda_1}=\frac{\sqrt{2}}{2}\begin{pmatrix*}[r] 1\\-1\end{pmatrix*}.
      \]
\item For $\lambda_2$ we have $$\begin{pmatrix*}[r]
  1&-1\\-1&1\end{pmatrix*} \begin{pmatrix*}[r] x_1
    \\ x_2\end{pmatrix*} = \begin{pmatrix*}[r] 0\\0\end{pmatrix*},$$
      which implies that the normalized eigenvector is
      \[
      x_{\lambda_2}=\frac{\sqrt{2}}{2}\begin{pmatrix*}[r]
        1\\1\end{pmatrix*}.
        \]
\end{itemize}

Thus, we can write $$Q= U \Lambda U^T = \begin{pmatrix*}[r] \frac{\sqrt{2}}{2}&\frac{\sqrt{2}}{2}\\-\frac{\sqrt{2}}{2}&\frac{\sqrt{2}}{2}\end{pmatrix*} 
\begin{pmatrix*}[r] 2&0\\0&0\end{pmatrix*} \begin{pmatrix*}[r] \frac{\sqrt{2}}{2}&\frac{\sqrt{2}}{2}\\-\frac{\sqrt{2}}{2}&\frac{\sqrt{2}}{2}\end{pmatrix*}^T.$$

  We have
  \[
  U^Tg= \frac{\sqrt{2}}{2} \begin{pmatrix*}[r]
    3\\3 \end{pmatrix*}.
  \]
As the second entry of this vector, corresponding to the zero eigenvalue, is not
zero, we can conclude that the problem is unbounded.

